\documentclass[11pt,a4paper]{article}
\usepackage[margin=1in]{geometry}
\usepackage{amsmath,amssymb,amsthm,booktabs,array}
\usepackage[colorlinks=true,allcolors=blue,breaklinks=true]{hyperref}
\newtheorem{theorem}{Theorem}
\newtheorem{lemma}[theorem]{Lemma}
\newtheorem{corollary}[theorem]{Corollary}
\newcommand{\R}{\mathbb R}
\newcommand{\C}{\mathbb C}
\newcommand{\conv}{\operatorname{conv}}
\newcommand{\cone}{\operatorname{cone}}
\newcommand{\diag}{\operatorname{diag}}
\title{A five-vertex negative answer to Johnson and Paparella's\protect\\
Question 9.1 on Hadamard-closed spectral polyhedra}
\author{Haobo Ma\\{\small ChronoAI Pte Ltd, Singapore}
\and Wenlin Zhang\\{\small National University of Singapore}\\
{\small ChronoAI Pte Ltd, Singapore}}
\date{November 13, 2026}
\begin{document}
\maketitle
\begin{abstract}
Question 9.1 of Johnson and Paparella, \emph{Perron similarities and the
nonnegative inverse eigenvalue problem}, arXiv:2409.07682v4, Section~9,
p.~27, has a negative answer for both polyhedral cones and polytopes.
For every integer $m\geq1$ we construct a rational polygon with $m+4$ edges
whose lift to $x_0=1$ and conical hull are closed under coordinatewise
multiplication and contain the all-ones vector. Every vector in the polygonal
lift is the spectrum, in a fixed coordinate order, of a symmetric stochastic
matrix with a common eigenbasis. Nevertheless neither set is the complete
spectratope or spectracone of any complex $3\times3$ matrix: such an exact
representation would constrain the facets away from the identity by its
three diagonal entries and their trace. The five-vertex case already fails
these constraints. Cartesian products give full-dimensional counterexamples
in every coordinate size $n\geq3$. For the complex-Euclidean boundary
used in that source, its Conjecture~9.2 also fails: the stochastic spectrum
$(1,0,\ldots,0)$ is a boundary point but is not radially extremal.
\end{abstract}

\section{The question and the construction}
For an invertible $n\times n$ complex matrix $S$, put
\begin{align*}
 M_S(x)&=S\diag(x)S^{-1},&
 \mathcal C(S)&=\{x\in\C^n:M_S(x)\geq0\},\\
 \mathcal P(S)&=\{x\in\mathcal C(S):M_S(x)e=e\},
\end{align*}
where $e=(1,\ldots,1)^T$, and $\geq0$ means real and entrywise
nonnegative. Johnson and Paparella's Question~9.1 asks whether every
polyhedral cone closed under the Hadamard product $x\circ y=(x_i y_i)_i$
is a set $\mathcal C(S)$, and likewise whether every Hadamard-closed
polytope is a set $\mathcal P(S)$ \cite[p.~27]{JP25}.
The matrix size is the number of coordinates, and equality of sets, rather
than containment, is required. We give counterexamples to both statements,
even with an identity element and a common stochastic realization.
The question begins: ``If $K$ is a polyhedral cone (polytope) that is closed
with respect to the Hadamard product'', and asks for an invertible $S$
giving the respective equality.

The real theory of Perron spectratopes was developed in \cite{JP16}; row-cone
criteria appear in \cite{JP17}. The complex definitions above and their
Hadamard-closure properties are Definition~4.1 and Theorem~4.4 of
\cite{JP25}. The characterization of totally extremal ideal Perron
similarities by character tables in \cite[Corollary~3.4]{AP26} concerns
matrices already satisfying those additional hypotheses, rather than arbitrary
Hadamard-closed polyhedra. Character tables provide ideal similarities
under the group-theoretic hypotheses in \cite{GLP26}.

Fix a positive integer $m$. For $0\leq j\leq m$ define
\[
 q_j=\left(\frac{j^2-m^2}{4(j^2+m^2)},
             -\frac{mj}{2(j^2+m^2)}\right),\qquad
 r=(1/4,0),\quad t=(0,1/4),\quad h=(1,1).
\]
Write
\[
 B_m=\{q_0,\ldots,q_m,r,t\},\qquad Q_m=\conv(B_m\cup\{h\}),
\]
\[
 P_m=\{(1,u,v):(u,v)\in Q_m\},\qquad K_m=\cone(P_m).
\]
Here $\conv$ denotes convex hull and $\cone$ denotes all finite nonnegative
linear combinations, including the empty combination $0$.
We regard $P_m,K_m\subset\R^3$ also as subsets of $\C^3$.

\begin{theorem}\label{thm:main}
For every integer $m\geq1$, the following hold.
\begin{enumerate}
\item $P_m$ is a rational two-dimensional polytope with $m+4$ edges.
The pointed cone $K_m$ is full-dimensional in $\R^3$ and has $m+4$ facets.
Both sets are polyhedra in the real-half-space convention on $\C^3$.
\item Both sets contain $e=(1,1,1)$ and are closed under the Hadamard product.
\item There is a real invertible matrix $S_0$ with
$P_m\subset\mathcal P(S_0)$ and $K_m\subset\mathcal C(S_0)$; the
realizing matrices can all be chosen symmetric.
\item For every $S\in\operatorname{GL}_3(\C)$,
$P_m\ne\mathcal P(S)$ and $K_m\ne\mathcal C(S)$.
\end{enumerate}
\end{theorem}

\section{Geometry and multiplication}
Every $q_j$ lies on $u^2+v^2=1/16$ in the closed negative quadrant, and
$q_0=(-1/4,0)$, $q_m=(0,-1/4)$. In particular, $Q_m$ contains the diamond
\begin{equation}\label{eq:diamond}
 D=\conv\{(1/4,0),(0,1/4),(-1/4,0),(0,-1/4)\}
   =\{(u,v):|u|+|v|\leq1/4\}.
\end{equation}
Thus it has nonempty interior. All points in $B_m$ belong to
$[-1/4,1/4]^2$.

For $0\leq k<m$ put $d_k=m^2+k(k+1)$ and
\[
 a_k=\frac{4(m^2-k(k+1))}{d_k},\qquad
 b_k=\frac{4m(2k+1)}{d_k},\qquad
 \ell_k(u,v)=1+a_k u+b_k v.
\]
Both coefficients are positive. Direct substitution gives the identity
\begin{equation}\label{eq:slack}
 \ell_k(q_j)=
 \frac{2m^2(j-k)(j-k-1)}{d_k(j^2+m^2)}\geq0.
\end{equation}
Since $j$ is integral, equality occurs exactly for $j=k,k+1$.
Also $\ell_k(r),\ell_k(t),\ell_k(h)>0$.
Consequently $[q_k,q_{k+1}]$ is an exposed edge of $Q_m$.
The four remaining edges, in cyclic order, have the following supporting
forms, nonnegative on $Q_m$:
\[
\begin{array}{c|c}
\text{edge}&\text{form}\\ \hline
{[q_m,r]}&1-4u+4v\\
{[r,h]}&1-4u+3v\\
{[h,t]}&1+3u-4v\\
{[t,q_0]}&1+4u-4v
\end{array}
\]
For the second and third forms, their values at every $q_j$ are at least
$1/4$, since $u,v\in[-1/4,0]$. For the first form the value is nonnegative,
and vanishes among the $q_j$ only at $q_m$; for the fourth it vanishes only
at $q_0$. Substitution at $r,t,h$ gives exactly the displayed endpoint
equalities. This verifies each asserted supporting edge.
The $m+4$ distinct generators are therefore all vertices, in the cyclic order
$q_0,\ldots,q_m,r,h,t$, and these are all the edges.

The intersection of the corresponding supporting half-planes is $Q_m$:
the listed segments form its closed convex boundary, with every successive
pair joined at the indicated vertex. Homogenizing their forms gives $K_m$.
For completeness, these homogeneous inequalities force $x_0\geq0$ by adding
the first and fourth forms in the table. If $x_0=0$, those two forms give
$u=v$, and the second gives $u\leq0$; any $\ell_k$ then gives $u\geq0$.
Thus the only such point is $0$. If $x_0>0$, division by $x_0$ recovers
$Q_m$. Its edges correspond bijectively to the facets of $K_m$.
The diamond in \eqref{eq:diamond} shows that the cone spans $\R^3$;
its positive first coordinate away from zero shows pointedness.
The equations $\operatorname{Im}x_i=0$, and $x_0=1$ for $P_m$, can each
be imposed by two real half-spaces. This also verifies the source's
complex-Euclidean polyhedron convention.

If $p,q\in B_m$, both coordinates of $p\circ q$ have absolute value
at most $1/16$. Hence $p\circ q\in D\subset Q_m$.
Since $h\circ p=p$ and $h\circ h=h$, products of all generators of
$Q_m$ lie in $Q_m$. For convex combinations, bilinearity gives
\[
 \left(\sum_i\alpha_i p_i\right)\circ
 \left(\sum_j\beta_j p_j\right)
   =\sum_{i,j}\alpha_i\beta_j(p_i\circ p_j),
 \qquad \sum_{i,j}\alpha_i\beta_j=1.
\]
Thus $Q_m$, and therefore $P_m$, is Hadamard-closed. The same identity
with arbitrary nonnegative coefficients proves closure of $K_m$.
The lift of $h$ is the identity $e$ in both sets.

\section{A common stochastic realization}
Let
\[
 e=(1,1,1)^T,\quad f=(1,-1,0)^T,\quad g=(1,1,-2)^T,
 \qquad S_0=[\,e\ f\ g\,].
\]
These vectors are orthogonal, with squared norms $3,2,6$, and
$\det S_0=6$. Therefore
\begin{equation}\label{eq:matrix}
 A(u,v)=S_0\diag(1,u,v)S_0^{-1}
       =\frac{ee^T}{3}+u\frac{ff^T}{2}+v\frac{gg^T}{6}.
\end{equation}
This matrix is symmetric and $A(u,v)e=e$. Its potentially smallest entries
on $[-1/4,1/4]^2$ are bounded by
\[
 \frac13\mathbin{\pm}\frac u2+\frac v6\geq\frac16,\qquad
 \frac13-\frac v3\geq\frac14,\qquad
 \frac13+\frac{2v}{3}\geq\frac16.
\]
Hence $A(p)\geq0$ for $p\in B_m$, while $A(h)=I$.
The affine dependence in \eqref{eq:matrix} gives $A(u,v)\geq0$ throughout
$Q_m$. Thus $P_m\subset\mathcal P(S_0)$, and taking nonnegative multiples
gives $K_m\subset\mathcal C(S_0)$.
In particular, the obstruction below is to an exact representation, not
to the realizability of the individual spectra.

\section{The identity-facet obstruction}
We use the real-half-space interpretation of complex polyhedra; see
\cite{Klee59} and \cite[Section~3.1]{JP25} for the underlying convex geometry.
\begin{lemma}\label{lem:facets}
A full-dimensional polyhedron in a real affine space, defined there by
$q$ affine inequalities, has at most $q$ facets.
\end{lemma}
\begin{proof}
Discard forms identically zero on the affine space. At a relative interior
point of a facet, at least one remaining inequality is active, since strict
inequalities describe an interior point. A nonnegative affine form vanishing
at a relative interior point of a facet vanishes on the entire facet:
each sufficiently short segment in that facet through the point has both
endpoints nonnegative, so its two directional derivatives must be zero.
Its zero hyperplane therefore supports that facet. A single supporting
hyperplane cuts out a single face and cannot support two distinct facets.
Assigning one active form to each facet proves the bound.
\end{proof}

\begin{lemma}\label{lem:identity}
Let $T$ be a full-dimensional polytope in a real affine space $H\subset\C^n$
containing $e$. Suppose that, on $H$, it is defined by the real inequalities
$L_{ij}(x)\geq0$, where $L_{ij}(x)=(S\diag(x)S^{-1})_{ij}$ for
some $S\in\operatorname{GL}_n(\C)$.
At most $n$ facets of $T$ omit $e$. If there are exactly $n$, their
supporting affine forms, each normalized to have value $1$ at $e$,
sum to $\sum_kx_k$ on $H$.
\end{lemma}
\begin{proof}
Since $S\diag(e)S^{-1}=I$, an off-diagonal form vanishes at $e$.
The supporting face it defines, when nontrivial, therefore contains $e$.
At a relative interior point of a facet omitting $e$, an active nonzero
form exists by the proof of Lemma~\ref{lem:facets}; that form must be
diagonal. Each of the $n$ diagonal forms supports at most one facet.
If all $n$ are needed, each equals the normalized affine form of its
facet, because $L_{ii}(e)=1$. Finally,
$\sum_iL_{ii}(x)=\operatorname{tr}(S\diag(x)S^{-1})=\sum_kx_k$.
\end{proof}

\begin{proof}[Completion of Theorem~\ref{thm:main}]
For $S\in\operatorname{GL}_3(\C)$ write the nine complex linear forms
\[
 L_{ij}(x)=\bigl(S\diag(x)S^{-1}\bigr)_{ij}
          =\sum_{k=0}^2 S_{ik}(S^{-1})_{kj}x_k
 \qquad (0\leq i,j\leq2).
\]
Work on $H=\{(1,u,v):u,v\in\R\}$. If $K_m=\mathcal C(S)$,
its section by $H$ is $P_m$. If $P_m=\mathcal P(S)$, use that
equality directly. In either case the imaginary parts of all entries
vanish on the relative interior of $P_m$ and hence on all of $H$.
In the spectratope case the row-sum equations also vanish there and hence
on $H$. Consequently either proposed representation gives
\[
 P_m=\{x\in H:L_{ij}(x)\geq0
                              \text{ for all }i,j\}.
\]
Exactly two edges contain $e$, namely the lifts of $[r,h]$ and $[h,t]$.
Thus $m+2$ edges omit $e$. For $m\geq2$ this exceeds three,
contradicting Lemma~\ref{lem:identity}.

For $m=1$, the polygon is the pentagon
\[
 Q_1=\conv\{(-1/4,0),(0,-1/4),(1/4,0),(1,1),(0,1/4)\}.
\]
The three edge forms away from $h$, normalized at $h$, are
\[
 \frac{1+4u+4v}{9},\qquad 1-4u+4v,\qquad 1+4u-4v.
\]
At $(u,v)=(0,0)$ their sum is $19/9$, whereas the trace of a
matrix with ordered eigenvalues $(1,0,0)$ is $1$.
This contradicts the last assertion of Lemma~\ref{lem:identity}.
\end{proof}

\begin{corollary}\label{cor:alln}
For every $n\geq3$ and $m\geq1$, define
\[
 P_{m,n}=\{(1,u,v,z_1,\ldots,z_{n-3}):(u,v)\in Q_m,
                        \ 0\leq z_i\leq1\},\qquad
 K_{m,n}=\cone(P_{m,n}).
\]
These rational unital Hadamard-closed polyhedra have dimensions $n-1$
and $n$, respectively, and each has $m+2n-2$ relative facets.
Every point has a nonnegative realization in a common eigenbasis;
the realizations for $P_{m,n}$ are stochastic. Nevertheless
$P_{m,n}\ne\mathcal P(S)$ and $K_{m,n}\ne\mathcal C(S)$ for every
$S\in\operatorname{GL}_n(\C)$.
\end{corollary}
\begin{proof}
The product with the $(n-3)$-cube preserves Hadamard closure and adds
$2(n-3)$ facets. Its conical hull is pointed, and its facets correspond
to the product's facets because its first-coordinate section is bounded.
Put $Z=\diag(z_1,\ldots,z_{n-3})$ and let $b$ be the column vector
with entries $1-z_i$. The matrix
\[
 B(u,v,z)=\begin{pmatrix}
 A(u,v)&0\\ b e^T/3&Z
 \end{pmatrix}
\]
is nonnegative and stochastic. A common eigenbasis is
$e_n,(f,0),(g,0)$ followed by the last $n-3$ coordinate vectors;
the corresponding eigenvalues are $1,u,v,z_1,\ldots,z_{n-3}$.
Here $e_n$ is the all-ones vector in $\R^n$, while $e,f,g$ in the
block formula are the three-coordinate vectors used above.
Scaling $B$ supplies the cone realizations.

As in the theorem, a proposed exact representation restricts on $x_0=1$
to its $n^2$ real entry inequalities. There are $m+2$ polygonal facets
away from $e_n$, together with $n-3$ facets $z_i=0$.
If $m\geq2$, their number $m+n-1$ exceeds $n$.
For $m=1$, their normalized forms are the three pentagonal forms above
and the $z_i$. Their sum at $(1,0,\ldots,0)$ is again $19/9$, not $1$.
Lemma~\ref{lem:identity} excludes both exact representations.
\end{proof}

\begin{corollary}\label{cor:rhc}
Let $V_m$ be the $(m+4)\times3$ matrix whose rows are the vertices of
$P_m$. Then $V_m$ has rank three and is row Hadamard conic in the sense
of \cite[Definition~2.3]{JP17}: each product of two rows lies in its
row cone. Its row cone is nevertheless not $\mathcal C(S)$ for any
$S\in\operatorname{GL}_3(\C)$.
\end{corollary}
\begin{proof}
The row cone is $K_m$, so closure, rank, and nonrepresentation follow
from Theorem~\ref{thm:main}. This rectangular statement leaves intact the
square-invertible hypotheses of the row-cone criteria in \cite{JP17}.
\end{proof}

\begin{corollary}\label{cor:boundary}
Use the source's notation $\mathrm{SL}_1^n$ for ordered spectra of
stochastic $n\times n$ matrices with first coordinate $1$, and $E_n$ for
its radially extremal points: $x\in E_n$ means
$(1,\alpha x_1,\ldots,\alpha x_{n-1})\notin\mathrm{SL}_1^n$ for all
$\alpha>1$.
With boundary taken in the complex affine space
$H_n^{\C}=\{x\in\C^n:x_0=1\}$, or in $\C^n$ itself,
\[
 (1,0,\ldots,0)\in\partial\mathrm{SL}_1^n\setminus E_n
 \qquad(n\geq2).
\]
In particular, Conjecture~9.2, p.~27, fails under this ambient-boundary
convention for every $n>3$.
\end{corollary}
\begin{proof}
The matrix $e_ne_n^T/n$ is stochastic with ordered spectrum
$x=(1,0,\ldots,0)$.
Every real matrix has real trace, so $\mathrm{SL}_1^n$ lies in the
proper real hyperplane $\operatorname{Im}\sum_kx_k=0$ of $H_n^{\C}$.
Every neighborhood of $x$ in $H_n^{\C}$ contains
$(1,i\varepsilon,0,\ldots,0)$ for some nonzero real $\varepsilon$;
its nonreal trace excludes it from $\mathrm{SL}_1^n$.
Thus $x$ is on the boundary. Multiplying its zero tail by any $\alpha>1$
leaves $x$ unchanged, so it is not extremal. The same perturbations
also establish the boundary assertion in $\C^n$.
\end{proof}
The boundary convention matters: the source uses complex-Euclidean
neighborhoods in its proof of Theorem~3.17, pp.~9--10, and describes
Conjecture~9.2 as asserting that ``points on the boundary are extremal''.
The corollary concerns that convention, not an alternative boundary
relative to a space of conjugation-symmetric lists or to a real spectral
parameterization.

\section{Coverage and verification}
The same proof permits arbitrarily many facets in three coordinates. It
already gives counterexamples at $m=1$; at $m=10$ the construction has
fourteen facets and agrees with the rational cone in the following Lean source.
The extra listed vector $(1,0,0)$ is the average of the four lifted axis
vertices, so it does not change the cone.
The declaration \texttt{result} in the
\href{https://github.com/the-omega-institute/trureturing/blob/e277c80d045570daf43bd03b4e313f7ce457b4e2/D5/S0/FiniteGeometry/Polytopes/JohnsonPaparellaHadamardPolyhedraRefutation.lean}{pinned Lean source}
proves the negation of \texttt{claim}, the joint universal converse of
Question~9.1, using the $m=10$ cone. Its cone decomposition, product closure,
and ten-distinct-supporting-forms obstruction concern that fixed cone.
The parameter family, exact facet counts, common stochastic realizations,
pentagonal trace obstruction, polytope conclusion, and
Corollaries~\ref{cor:alln}--\ref{cor:boundary} have the written proofs above.
The accompanying \texttt{verify.py} gives exact rational finite checks of
the displayed generators, supporting forms, convex-product certificates,
and stochastic matrix identities, including the pentagon, its cube products,
and the fixed $m=10$ example.

\begin{center}
\begin{tabular}{@{}>{\raggedright\arraybackslash}p{.36\linewidth}>{\raggedright\arraybackslash}p{.56\linewidth}@{}}
\toprule
Source item&Coverage\\
\midrule
Question 9.1, cone part&Answered: no; $K_m$, every $m\geq1$;
also $K_{m,n}$ for every $n\geq3$.\\
Question 9.1, polytope part&Answered: no; $P_m$, every $m\geq1$;
also $P_{m,n}$ for every $n\geq3$.\\
Theorems 4.4 and 4.6, forward closure and polyhedral properties&
Prior results cited; the examples have these properties, not their converses.\\
Definition 4.18 and Theorem 4.22, ideal similarities&
Prior characterization cited; containment in a spectratope does not imply equality.\\
Conjecture 9.2, boundary extremality for $n>3$&Refuted (explicit counterexample)
for the complex ambient boundary; Corollary~\ref{cor:boundary}.\\
All-size and five-vertex variants&Not posed separately; proved above.\\
\cite[Definition 2.3]{JP17}, rectangular row Hadamard closure&
Not posed as a converse; obstruction proved in Corollary~\ref{cor:rhc}.\\
\cite[Remark 3.4]{JP17}, row realization of every nonnegative spectrum&
Not covered; no conclusion about that realization question.\\
\cite[Conjecture 7.1]{GLP26}, totally extremal ideal similarities&
Prior proof cited: Artemis--Paparella (2026), Corollary 3.4; not reproved here.\\
\bottomrule
\end{tabular}
\end{center}
The boundary counterexample uses the full stochastic spectral region,
not the boundary of one of our prescribed polyhedra.
The earlier real spectratope and row-cone papers \cite{JP16,JP17} and
the characterization in \cite[Corollary~3.4]{AP26} concern fixed
similarities or additional ideality and extremality assumptions;
their statements are not converses for arbitrary Hadamard-closed polyhedra.

\begin{thebibliography}{9}
\bibitem{Klee59} V. Klee,
\emph{Some characterizations of convex polyhedra},
Acta Math. \textbf{102} (1959), 79--107,
\href{https://doi.org/10.1007/BF02559569}{doi:10.1007/BF02559569}.
\bibitem{JP16} C. R. Johnson and P. Paparella,
\emph{Perron spectratopes and the real nonnegative inverse eigenvalue problem},
Linear Algebra Appl. \textbf{493} (2016), 281--300.
\href{https://arxiv.org/abs/1508.07400v2}{arXiv:1508.07400v2}.
\bibitem{JP17} C. R. Johnson and P. Paparella,
\emph{Row cones, Perron similarities, and nonnegative spectra},
Linear Multilinear Algebra \textbf{65} (2017), 2124--2130.
\href{https://arxiv.org/abs/1611.02752v1}{arXiv:1611.02752v1}.
\bibitem{JP25} C. R. Johnson and P. Paparella,
\emph{Perron similarities and the nonnegative inverse eigenvalue problem},
Trans. Amer. Math. Soc. \textbf{378} (2025), 8361--8389,
\href{https://doi.org/10.1090/tran/9596}{doi:10.1090/tran/9596}.
Posing text: \href{https://arxiv.org/abs/2409.07682v4}{arXiv:2409.07682v4},
Definition 4.1, p.~10; Question 9.1 and Conjecture 9.2, p.~27.
\bibitem{AP26} E. J. Artemis and P. Paparella,
\emph{On a conjecture concerning totally extremal ideal Perron similarities},
\href{https://arxiv.org/abs/2606.02865v1}{arXiv:2606.02865v1} (2026).
\bibitem{GLP26} D. Z. Gershnik, A. J. Lewis and P. Paparella,
\emph{Character tables are ideal Perron similarities},
J. Algebra \textbf{694} (2026), 782--800,
\href{https://doi.org/10.1016/j.jalgebra.2026.02.004}{doi:10.1016/j.jalgebra.2026.02.004}.
\href{https://arxiv.org/abs/2508.02830v1}{arXiv:2508.02830v1} (2025),
Conjecture 7.1, p.~21.
\end{thebibliography}
\noindent\small\copyright\ 2026 Haobo Ma and Wenlin Zhang.
Licensed under \href{https://creativecommons.org/licenses/by/4.0/}{CC BY 4.0}.
\end{document}
